% Round 17. Foundation branch. Candidate; independent review required.
% Uses the separately stated C217 all-characteristic quadric criterion.
\section{Degenerate ordinary quadratic points at two and the all-prime reduction sector}
\label{r17cd2:section}

\paragraph{Status and published input.}
\textbf{Internally reviewed.} The final claims register identifies the independent audits and frozen dependencies.
This section treats the even-dimensional ordinary quadratic
case omitted from C217. Its local normal form is the published
Jannsen--Saito Lemma 4(ii), Section 3, printed p.11, and its
order reduction is their Theorem 5(d), printed pp.11--14,
\url{https://arxiv.org/pdf/0911.1470}. Their lemma explicitly
allows a finite etale extension of the DVR and gives an
isomorphism of henselized local rings. We first explain why
that statement supplies common pointed etale neighborhoods.
We then check the native component geometry needed for the
all-characteristic numerical quadric-forest theorem C217.
No claim of novelty is made for the normal forms or blowups.

\subsection{Using the published henselian normal form in actual etale charts}

\begin{lemma}[A henselian local isomorphism gives a common etale neighborhood]
\label{r17cd2:henselian-chart}
Let $A$ and $B$ be local rings of finite-type schemes over
a noetherian ring $R$, with the same residue field at the
chosen points. Suppose there is an isomorphism of local
$R$-algebras $A^h\simeq B^h$ inducing that residue
identification. Then the two schemes have a common pointed
etale neighborhood with that residue field.
\end{lemma}

\begin{proof}
Choose an affine finite-presentation algebra representing
a neighborhood on the $A$ side. Its map into $B^h$ factors
through one of the finite-presentation pointed etale
neighborhoods defining that henselization. Localizing at
the selected point gives a local $R$-map $A\to B'$, with
$B'$ pointed etale over $B$ and with unchanged residue
field. Its induced map $A^h\to(B')^h\simeq B^h$ is the
given isomorphism by the universal property of henselization.

The induced map of completed noetherian local rings is
therefore an isomorphism. Since $\widehat {B'}$ is
faithfully flat over $B'$ and is identified over $A$ with
the flat completion $\widehat A$, flatness descends to
$A\to B'$. The completed tangent-space equality and
the unchanged residue field make
$\Omega_{B'/A}\otimes\kappa(B')=0$. This module is
finite, so local Nakayama gives $\Omega_{B'/A}=0$.
Thus the essentially finite-presentation map is flat and
unramified, hence etale at the point. Shrink the selected
neighborhood so that both maps are etale. This proves the
claim by finite-presentation descent; it does not attempt
to algebraize an arbitrary formal isomorphism.
\end{proof}

Let $K/\mathbf Q_2$ have valuation ring $R$, and let
$\mathcal X/R$ be integral, flat and projective with
smooth geometrically integral generic fiber of even
dimension $n=2a\ge4$. Assume its geometrically integral
special fiber is smooth except at finitely many ordinary
quadratic points. Here ordinary quadratic means that
the projectivized geometric tangent cone is a smooth
quadric. In this even relative dimension its quadratic
form has a one-dimensional polar radical in residue two;
it is the degenerate case in the terminology of
Jannsen--Saito.

\begin{proposition}[A common finite field for all degenerate node equations]
\label{r17cd2:field}
There is a finite extension $L/K$ such that, at every
special node of $\mathcal X_{R_L}$, there is a common
pointed etale neighborhood with a hypersurface
\[
  P(x_1,\ldots,x_n)+z^2+\epsilon\varpi_L^qz=0,
       \qquad q\ge1,\quad \epsilon\in R_L^\times,
 \tag{CD217.1}\label{r17cd2:split-polynomial}
\]
whose polar form on the $n$ variables of $P$ is unimodular.
One may arrange that the selected closed nodes are rational
over the residue field of $L$. There is no assertion that
$[L:K]$ is two or is independent of the number of nodes.
\end{proposition}

\begin{proof}
First take a common finite unramified extension $K_0/K$
which contains the finitely many residue fields and the
finite etale DVR extensions allowed in the published
normal-form lemma. Its conjugates give the same forms
at all conjugate nodes; enlarging it finitely if needed
handles all of them at once. The nodes are rational over
its residue field. For each one, Jannsen--Saito Lemma 4(ii)
gives a henselian local equation
\[
                   P(x)+z^2+bz+c=0,
       \qquad 0\ne\Delta=b^2-4c\in\mathfrak m_{R_0},
 \tag{CD217.2}\label{r17cd2:js-form}
\]
with $P$ nondegenerate in its even number $n$ of variables.
The chosen special point is the origin, so $c$ lies in
the maximal ideal; also $b$ does, since $b^2-4c$ does
and the residue characteristic is two. Apply
Lemma~\ref{r17cd2:henselian-chart} to obtain actual common
pointed etale charts.

Take a finite extension $L/K_0$ splitting all the finitely
many polynomials $T^2+bT+c$, including the conjugate data.
For example one may take a normal closure of the compositum
of their quadratic splitting fields over $K$. Retain the
full resulting degree. A root $\alpha\in L$ is integral,
being a root of a monic polynomial, and its reduction is
zero because the reduced polynomial is $T^2$. Thus the
translation $z=w+\alpha$ fixes the special origin and
changes \eqref{r17cd2:js-form} to
\[
                    P(x)+w^2+\beta w=0,
                    \qquad \beta=2\alpha+b.
\]
One has $\beta^2=\Delta\ne0$. Its valuation is positive,
so $\beta=\epsilon\varpi_L^q$ for a finite integer
$q\ge1$. The polar matrix of $P$ stays invertible after
the extension. Renaming $w$ as $z$ proves the assertion.
All chosen charts remain etale after base change.

This finite extension is auxiliary. In the final detector
its complete degree $[L:K]$, including residue extension
and ramification, will be retained. No selected geometric
node or isolated square-root branch replaces that degree.
\end{proof}

\subsection{The order-one reduction and its exact exceptional quadrics}

\begin{proposition}[Native degenerate-node chains]
\label{r17cd2:chains}
For the local equation \eqref{r17cd2:split-polynomial},
exactly $q$ successive blowups at the remaining reduced
special point produce a regular reduced strictly semistable
model. Its exceptional chain consists of $q-1$ copies of
\[
       \mathbf P_E(\mathcal O_E\oplus\mathcal O_E(-1))
\]
and a terminal smooth projective quadric $T$, where
\[
 \begin{split}
 E&=\{\overline P(X)+Z^2=0\}\subset\mathbf P^n_{k_L},\\
 T&=\{\overline P(X)+Z^2+\overline\epsilon ZS=0\}
                                      \subset\mathbf P^{n+1}_{k_L}.
 \end{split}
 \tag{CD217.3}\label{r17cd2:quadrics}
\]
Every double stratum is a copy of the smooth odd-dimensional
quadric $E$, every component has multiplicity one, and the
two sections on an intermediate component are disjoint.
These statements remain valid over the native residue field
when the selected node has not been made rational.
\end{proposition}

\begin{proof}
Write $\pi=\varpi_L$ for this calculation. Blow up the
maximal ideal $(x_1,\ldots,x_n,z,\pi)$.
The saturated $x_i$-chart, with $t=x_i$, $z=tw$ and
$\pi=tv$, is
\[
   tv=\pi,\qquad
   P(y_1,\ldots,1,\ldots,y_n)+w^2
                     +\epsilon t^{q-1}v^q w=0.
 \tag{CD217.4}\label{r17cd2:x-chart}
\]
The saturated $z$-chart, with $t=z$, is
\[
          tv=\pi,\qquad
          P(y)+1+\epsilon t^{q-1}v^q=0.
 \tag{CD217.5}\label{r17cd2:z-chart}
\]
The saturated $\pi$-chart, with $x_i=\pi u_i$ and
$z=\pi w$, is
\[
                  P(u)+w^2+\epsilon\pi^{q-1}w=0.
 \tag{CD217.6}\label{r17cd2:pi-chart}
\]
These are actual Rees charts. In the first two charts
the remaining polynomial modulo $t$ is a nonzero
polynomial in the domain $k_L[v,y,w]$ or $k_L[v,y]$.
The regular-sequence permutation argument in the
semistable ambient $tv=\pi$ makes multiplication by $t$
injective after imposing it, so removal of $t^2$ gives
the exact saturation. In the last chart the equation is
monic in $w$, hence its quotient is flat over the DVR
and has no $\pi$-torsion. This gives its exact saturation.

Let $B_P$ be the invertible polar matrix of $\overline P$.
The quadric $E$ is smooth: vanishing of its $X$ derivatives
forces $X=0$, and its equation then forces $Z=0$,
impossible projectively. The same argument shows that
the cone $\overline P(u)+w^2=0$ has the origin as its
only reduced singular point.

For $q>1$, the last summands of
\eqref{r17cd2:x-chart} and \eqref{r17cd2:z-chart} vanish
on the reduced special fiber $tv=0$. On an $x_i$-chart
some $y$ derivative of $\overline P(1,y)$ is nonzero:
if all the other $X$ derivatives vanished, Euler's
identity in characteristic two would force the $i$
derivative to vanish too, contradicting $B_PX\ne0$.
On a $z$-chart, $\overline P(y)+1=0$ implies $y\ne0$,
so some $y$ derivative is nonzero. Etale elimination
therefore leaves the regular semistable equation
$tv=\pi$. In the $\pi$-chart the only remaining special
singular point is the origin, and the coefficient order
has decreased from $q$ to $q-1$.

For $q=1$, the $\pi$-chart is smooth because its derivative
in $w$ reduces to $\overline\epsilon\ne0$.
In the $x_i$- and $z$-charts, at the double locus
$t=v=0$ the preceding $y$-derivative argument still
applies and leaves $tv=\pi$. Where $v$ is a unit,
one solves $t=\pi/v$; in the $x_i$-chart the $w$
derivative is then a unit, while in the $z$-chart the
coefficient of $v$ is a unit. These are smooth
one-component charts. Where $t$ is a unit the same
$y$-derivative test gives a smooth chart after solving
$v=\pi/t$. Thus every chart is regular and strictly
semistable at the last blowup.

At $q>1$ the exceptional component is the projective
cone over $E$, with vertex away from its preceding
double stratum. At $q=1$ it is the terminal quadric $T$
in \eqref{r17cd2:quadrics}. Its polar matrix is the
direct sum of $B_P$ and
\[
             \begin{pmatrix}0&\overline\epsilon\\
                             \overline\epsilon&0\end{pmatrix},
\]
so it is invertible and $T$ is smooth. Its hyperplane
section $S=0$ is $E$.

Blowing up the cone vertex gives
$\mathbf P_E(\mathcal O_E\oplus\mathcal O_E(-1))$,
by the direction-chart construction used in QF17 and
C217. Its vertex and infinity sections are disjoint.
Induction now gives exactly $q$ blowups, $q-1$
intermediate bundles and one terminal quadric. The
structural equations $tv=\pi$ give multiplicity one,
and the moving singular vertex is disjoint from every
previous double stratum, so no triple stratum occurs.
Pointed etale base change identifies the native graded
cones and the corresponding exceptional schemes; no
splitting of their ruling classes is needed.
\end{proof}

\subsection{All ordinary quadratic reductions in mixed characteristic}

\begin{theorem}[All-degree actual G for even-dimensional ordinary quadratic reduction at two]
\label{r17cd2:actual}
Let $K/\mathbf Q_2$ and its absolute value be fixed,
and let $X/K$ be smooth projective geometrically integral
of even dimension $n\ge4$. Suppose it has an integral
flat projective model whose geometrically integral special
fiber is smooth except at finitely many ordinary quadratic
points. Then every geometrically homologically trivial
rational cycle in every positive codimension has strong
ordinary actual G over $K$, on its exact specified support
complement and against every compact original actual test.
The same conclusion holds for targets admitting such a
projective dominant witness after a finite auxiliary
extension, and anew for all cycles after every finite
extension of $K$.
\end{theorem}

\begin{proof}
Use Proposition~\ref{r17cd2:field} once for all nodes.
The base-changed model remains integral because it is
flat with geometrically integral generic fiber.
The special geometric nodes persist. Apply the preceding
native blowup construction at each one. Away from them
flatness and fiber smoothness make the morphism smooth.
The resulting model is regular, projective and reduced
strictly semistable. Its main component and disjoint
exceptional chains form a tree, with all double strata
the native smooth quadrics $E$.

The all-characteristic quadric-forest theorem C217 gives
numerical admissible extensions in every codimension.
Here the double quadrics have odd dimension, so only
the one-dimensional ample numerical lines occur; no
middle-degree sign argument is needed. The frozen native
numerical witness and dominant Bott--Chern detectors
descend obstruction zero through the full finite extension
$L/K$, with its complete degree $[L:K]$ and all residue
weights. No degree-two shortcut is used for the simultaneous
splitting of several node polynomials. Reconstructing from
the curved pair for the original specified cycle gives
the asserted exact-support all-actual solution.

The dominant-target assertion is the same frozen detector.
After any finite extension the geometric ordinary
quadratic special points, flatness and smooth generic
fiber remain. Applying this theorem afresh gives the
full conclusion for cycles newly defined over that field.
No proper actual-form pushforward or off-support
finiteness of an alteration is asserted.
\end{proof}

\begin{corollary}[The ordinary quadratic local theorem at every prime]
\label{r17cd2:all-primes}
For every finite $K/\mathbf Q_p$, every dimension $n\ge3$,
and every smooth projective geometrically integral $X/K$
with an integral flat projective ordinary quadratic model
having geometrically integral special fiber, all
geometrically homologically trivial rational cycles satisfy
strong ordinary actual G in every positive codimension.
The model may have arbitrary node orders and need not be
regular. The conclusion is stable under finite field
extension and the stated projective dominant witnesses.
\end{corollary}

\begin{proof}
Use ON17 for odd residue characteristic, C217 for odd
dimension at residue two, and the preceding theorem for
even dimension at residue two. These cases exhaust the
stated range. Each conclusion retains the original field,
exact support and all actual tests.
\end{proof}

\begin{corollary}[Complete height on the all-prime ordinary quadratic reduction sector]
\label{r17cd2:height}
Let $X/F$ be smooth projective geometrically integral
over a number field. Suppose that at every finite place,
after a finite auxiliary local extension if needed, it
has a projective dominant witness with good reduction
or an ordinary quadratic model in the range of
Corollary~\ref{r17cd2:all-primes}.
Then the canonical real Beilinson--Bloch height is defined
on the complete geometric homological rational Chow
groups in every pair of complementary degrees, with the
frozen number-field normalization and correspondence
adjunction. GA17 gives its normalized admissible
arithmetic-intersection description on a fixed regular
arithmetic alteration.

In particular there is no residue-two exclusion for
varieties which themselves have good or ordinary quadratic
reduction at every finite place. This remains a specified
reduction sector, not a theorem that every variety has
such a reduction witness.
\end{corollary}

\begin{proof}
The local theorem and the good-reduction result give zero
obstruction for every homological class at every finite
place. Hence each complete homological Chow group equals
the corresponding frozen height domain. Apply that
canonical height theorem and GA17. The local auxiliary
extensions need not be completions of one prescribed
global extension: native local descent occurs first,
before the global height theorem is invoked.
\end{proof}
